% Options for packages loaded elsewhere
\PassOptionsToPackage{unicode}{hyperref}
\PassOptionsToPackage{hyphens}{url}
\PassOptionsToPackage{dvipsnames,svgnames,x11names}{xcolor}
%
\documentclass[
  10pt,
]{article}
\usepackage{amsmath,amssymb}
\usepackage{iftex}
\ifPDFTeX
  \usepackage[T1]{fontenc}
  \usepackage[utf8]{inputenc}
  \usepackage{textcomp} % provide euro and other symbols
\else % if luatex or xetex
  \usepackage{unicode-math} % this also loads fontspec
  \defaultfontfeatures{Scale=MatchLowercase}
  \defaultfontfeatures[\rmfamily]{Ligatures=TeX,Scale=1}
\fi
\usepackage{lmodern}
\ifPDFTeX\else
  % xetex/luatex font selection
  \setmainfont[]{DejaVu Serif}
  \setsansfont[]{DejaVu Sans}
  \setmonofont[]{DejaVu Sans Mono}
  \setmathfont[]{Latin Modern Math}
\fi
% Use upquote if available, for straight quotes in verbatim environments
\IfFileExists{upquote.sty}{\usepackage{upquote}}{}
\IfFileExists{microtype.sty}{% use microtype if available
  \usepackage[]{microtype}
  \UseMicrotypeSet[protrusion]{basicmath} % disable protrusion for tt fonts
}{}
\makeatletter
\@ifundefined{KOMAClassName}{% if non-KOMA class
  \IfFileExists{parskip.sty}{%
    \usepackage{parskip}
  }{% else
    \setlength{\parindent}{0pt}
    \setlength{\parskip}{6pt plus 2pt minus 1pt}}
}{% if KOMA class
  \KOMAoptions{parskip=half}}
\makeatother
\usepackage{xcolor}
\usepackage[margin=2.0cm]{geometry}
\usepackage{color}
\usepackage{fancyvrb}
\newcommand{\VerbBar}{|}
\newcommand{\VERB}{\Verb[commandchars=\\\{\}]}
\DefineVerbatimEnvironment{Highlighting}{Verbatim}{commandchars=\\\{\}}
% Add ',fontsize=\small' for more characters per line
\newenvironment{Shaded}{}{}
\newcommand{\AlertTok}[1]{\textcolor[rgb]{1.00,0.00,0.00}{\textbf{#1}}}
\newcommand{\AnnotationTok}[1]{\textcolor[rgb]{0.38,0.63,0.69}{\textbf{\textit{#1}}}}
\newcommand{\AttributeTok}[1]{\textcolor[rgb]{0.49,0.56,0.16}{#1}}
\newcommand{\BaseNTok}[1]{\textcolor[rgb]{0.25,0.63,0.44}{#1}}
\newcommand{\BuiltInTok}[1]{\textcolor[rgb]{0.00,0.50,0.00}{#1}}
\newcommand{\CharTok}[1]{\textcolor[rgb]{0.25,0.44,0.63}{#1}}
\newcommand{\CommentTok}[1]{\textcolor[rgb]{0.38,0.63,0.69}{\textit{#1}}}
\newcommand{\CommentVarTok}[1]{\textcolor[rgb]{0.38,0.63,0.69}{\textbf{\textit{#1}}}}
\newcommand{\ConstantTok}[1]{\textcolor[rgb]{0.53,0.00,0.00}{#1}}
\newcommand{\ControlFlowTok}[1]{\textcolor[rgb]{0.00,0.44,0.13}{\textbf{#1}}}
\newcommand{\DataTypeTok}[1]{\textcolor[rgb]{0.56,0.13,0.00}{#1}}
\newcommand{\DecValTok}[1]{\textcolor[rgb]{0.25,0.63,0.44}{#1}}
\newcommand{\DocumentationTok}[1]{\textcolor[rgb]{0.73,0.13,0.13}{\textit{#1}}}
\newcommand{\ErrorTok}[1]{\textcolor[rgb]{1.00,0.00,0.00}{\textbf{#1}}}
\newcommand{\ExtensionTok}[1]{#1}
\newcommand{\FloatTok}[1]{\textcolor[rgb]{0.25,0.63,0.44}{#1}}
\newcommand{\FunctionTok}[1]{\textcolor[rgb]{0.02,0.16,0.49}{#1}}
\newcommand{\ImportTok}[1]{\textcolor[rgb]{0.00,0.50,0.00}{\textbf{#1}}}
\newcommand{\InformationTok}[1]{\textcolor[rgb]{0.38,0.63,0.69}{\textbf{\textit{#1}}}}
\newcommand{\KeywordTok}[1]{\textcolor[rgb]{0.00,0.44,0.13}{\textbf{#1}}}
\newcommand{\NormalTok}[1]{#1}
\newcommand{\OperatorTok}[1]{\textcolor[rgb]{0.40,0.40,0.40}{#1}}
\newcommand{\OtherTok}[1]{\textcolor[rgb]{0.00,0.44,0.13}{#1}}
\newcommand{\PreprocessorTok}[1]{\textcolor[rgb]{0.74,0.48,0.00}{#1}}
\newcommand{\RegionMarkerTok}[1]{#1}
\newcommand{\SpecialCharTok}[1]{\textcolor[rgb]{0.25,0.44,0.63}{#1}}
\newcommand{\SpecialStringTok}[1]{\textcolor[rgb]{0.73,0.40,0.53}{#1}}
\newcommand{\StringTok}[1]{\textcolor[rgb]{0.25,0.44,0.63}{#1}}
\newcommand{\VariableTok}[1]{\textcolor[rgb]{0.10,0.09,0.49}{#1}}
\newcommand{\VerbatimStringTok}[1]{\textcolor[rgb]{0.25,0.44,0.63}{#1}}
\newcommand{\WarningTok}[1]{\textcolor[rgb]{0.38,0.63,0.69}{\textbf{\textit{#1}}}}
\usepackage{longtable,booktabs,array}
\usepackage{calc} % for calculating minipage widths
% Correct order of tables after \paragraph or \subparagraph
\usepackage{etoolbox}
\makeatletter
\patchcmd\longtable{\par}{\if@noskipsec\mbox{}\fi\par}{}{}
\makeatother
% Allow footnotes in longtable head/foot
\IfFileExists{footnotehyper.sty}{\usepackage{footnotehyper}}{\usepackage{footnote}}
\makesavenoteenv{longtable}
\setlength{\emergencystretch}{3em} % prevent overfull lines
\providecommand{\tightlist}{%
  \setlength{\itemsep}{0pt}\setlength{\parskip}{0pt}}
\setcounter{secnumdepth}{-\maxdimen} % remove section numbering
% doc_header.tex -- LaTeX preamble additions used by build_docs.sh for every
% Markdown -> LaTeX -> PDF conversion (the documents of the folder Revision_HS_CJT_and_Kohn-Sham).
\usepackage{etoolbox}
% longtable may not be loaded by pandoc when a document has no table
\usepackage{longtable}
% bold vectors: \bm -> \symbf under unicode-math (XeLaTeX), bm package otherwise
\makeatletter
\@ifpackageloaded{unicode-math}{\providecommand{\bm}[1]{\symbf{#1}}}{\usepackage{bm}}
\makeatother
% smaller type in long tables so that wide audit tables fit the text width;
% inside a table a line break is allowed after an underscore of plain text
% (pandoc writes it as \_), so that long identifiers wrap inside their column
% rows are set a little apart, so that symbols with indices do not touch
\AtBeginEnvironment{longtable}{\footnotesize\renewcommand{\arraystretch}{1.45}%
  \let\sptunderscore\_\def\_{\sptunderscore\allowbreak}}
\setlength{\LTcapwidth}{\textwidth}
% A short table is kept on one page.  The build (code/scripts/keep_tables.py) ends the rows of such a table with \\*,
% which forbids a page break after a row, and sets \sptkeptpenalty to 10000 in front of the table.  This penalty is
% put here in the place of the breakpoint (penalty 0) that longtable leaves after every table, so that the paragraph
% that follows a kept table, its caption, begins on the same page; it is set back to 0 after the table.
\newcount\sptkeptpenalty
\sptkeptpenalty=0
\makeatletter
\patchcmd{\endlongtable}{\endgraf\penalty\z@\addvspace\LTpost}%
  {\endgraf\penalty\sptkeptpenalty\addvspace\LTpost\global\sptkeptpenalty=\z@}%
  {}{\PackageError{doc_header}{the end of longtable could not be patched}{The tables can still be typeset, but a short table may be separated from its caption.}}
\makeatother
% allow line breaks after slashes and hyphens inside table cells
\usepackage{xurl}
% code blocks: break over-long lines instead of running into the margin.
% Plain ``` blocks are typeset by pandoc as {verbatim}, highlighted ones
% (```bash) as {Highlighting}; both are re-declared as breakable fancyvrb
% environments.  The continuation line is marked with a small hooked arrow.
\usepackage{fvextra}
\RecustomVerbatimEnvironment{verbatim}{Verbatim}{breaklines,breakanywhere,fontsize=\small}
\makeatletter
\@ifundefined{Highlighting}{}{%
  \RecustomVerbatimEnvironment{Highlighting}{Verbatim}{breaklines,breakanywhere,fontsize=\small,commandchars=\\\{\}}}
\makeatother
% inline code (pandoc writes it as \texttt{...}): never hyphenate -- an inserted
% hyphen could be read as part of a file name -- but allow a line break after
% an underscore, slash, comma, colon or equals sign, so that long paths and
% column lists wrap inside the text width and inside table columns.
% Inside a formula (e.g. the name of a notebook symbol in an equation) the same
% command must first leave math mode, which \text of amsmath does.
\ExplSyntaxOn
\tl_new:N \l__spt_code_tl
\cs_new_protected:Npn \__spt_texttt:n #1
  {
    \tl_set:Nn \l__spt_code_tl {#1}
    \regex_replace_all:nnN { ([\/\,\:\=]) } { \1 \c{allowbreak} } \l__spt_code_tl
    \regex_replace_all:nnN { \c{_} } { \c{_} \c{allowbreak} } \l__spt_code_tl
    \group_begin: \ttfamily \hyphenchar \font = -1 \scan_stop: \l__spt_code_tl \group_end:
  }
\RenewDocumentCommand \texttt { m }
  {
    \mode_if_math:TF { \text { \__spt_texttt:n {#1} } } { \__spt_texttt:n {#1} }
  }
\ExplSyntaxOff
% A figure is written in the Markdown source as an image paragraph followed by
% a caption paragraph ("**Figure S1.** ..."), so that the caption is also
% visible on GitHub.  An image that starts a paragraph is centered.
\usepackage{graphicx}% pandoc loads it only for documents that contain an image
\NewCommandCopy\sptincludegraphics\includegraphics
\RenewDocumentCommand\includegraphics{O{} m}{%
  \ifvmode\begin{center}\sptincludegraphics[#1]{#2}\end{center}%
  \else\sptincludegraphics[#1]{#2}\fi}
% keep figures where they are referenced
\usepackage{float}
\floatplacement{figure}{H}
% running header with the document title
\usepackage{fancyhdr}
\pagestyle{fancy}
\fancyhf{}
\fancyhead[L]{\small PROJECT spacetimes, Revision\_HS\_CJT\_and\_Kohn-Sham}
\fancyhead[R]{\small\thepage}
\renewcommand{\headrulewidth}{0.4pt}
% 16x16 matrices: tighter columns so that a full matrix fits the text width
\setlength{\arraycolsep}{3.2pt}
% The documents of this project contain long tagged symbols, file paths and
% identifiers that cannot be hyphenated.  A paragraph is therefore allowed to
% have looser lines before TeX gives up and lets a line stick out of the margin.
\AtBeginDocument{\tolerance=2000 \emergencystretch=5em \hbadness=2000}
\ifLuaTeX
  \usepackage{selnolig}  % disable illegal ligatures
\fi
\IfFileExists{bookmark.sty}{\usepackage{bookmark}}{\usepackage{hyperref}}
\IfFileExists{xurl.sty}{\usepackage{xurl}}{} % add URL line breaks if available
\urlstyle{same}
\hypersetup{
  pdftitle={Revision\_HS\_CJT\_and\_Kohn-Sham: the superfluid boson singlet state of the gauge-covariant spinor field on the 3+1 manifold of the Lorentz super-loop, by the methods of Hubbard-Stratonovich, Cornwall-Jackiw-Tomboulis and Kohn-Sham},
  pdfauthor={Prepared for Patrick Lee Nash, Ph.D.},
  colorlinks=true,
  linkcolor={blue},
  filecolor={Maroon},
  citecolor={blue},
  urlcolor={blue},
  pdfcreator={LaTeX via pandoc}}

\title{Revision\_HS\_CJT\_and\_Kohn-Sham: the superfluid boson singlet
state of the gauge-covariant spinor field on the 3+1 manifold of the
Lorentz super-loop, by the methods of Hubbard-Stratonovich,
Cornwall-Jackiw-Tomboulis and Kohn-Sham}
\usepackage{etoolbox}
\makeatletter
\providecommand{\subtitle}[1]{% add subtitle to \maketitle
  \apptocmd{\@title}{\par {\large #1 \par}}{}{}
}
\makeatother
\subtitle{README of the folder Revision\_HS\_CJT\_and\_Kohn-Sham of
PROJECT spacetimes}
\author{Prepared for Patrick Lee Nash, Ph.D.}
\date{November 21, 2026}

\begin{document}
\maketitle

\hypertarget{what-this-folder-is}{%
\section{1. What this folder is}\label{what-this-folder-is}}

This folder of PROJECT spacetimes (repository
\texttt{https://github.com/once-ere/spacetimes.git}) contains the
solution of the problem that the author posed on 20 November 2026: the
gauge-covariant field equations of a self-interacting spinor field on
the \(3+1\) manifold of the Lorentz super-loop, their analytic
solutions, and the condensation of the spinor field into a superfluid
boson singlet state of spin zero, described by three methods
(Hubbard-Stratonovich, Cornwall-Jackiw-Tomboulis and Kohn-Sham), with
the effective Lagrangian of each method, the spontaneous breaking of the
symmetry and the masses of the gauge bosons.

Every document of the folder is complete in itself: every equation,
symbol and definition that it uses is written in it. This README is the
list of contents of the folder.

\hypertarget{state}{%
\section{2. State}\label{state}}

\textbf{20 November 2026.} The folder has been created; the task record
and this README are written. The work paused after that, as the author's
Rule 2 requires after a false statement has been recorded with its cause
(I had stated that the reading of four symbol registries of the earlier
work was complete when two of them had not yet been read completely; the
statement was corrected in the next message). At 18:28 PST the author
instructed to continue and to complete all stages; at 19:25 PST, to push
every item as soon as it is complete. The work continues; this table is
updated with every item.

\textbf{Part I of the manuscript} (items 1 to 5 of the problem: the
manifold M31, the spinor field and theorem D1, the gauge group GEE and
its covariant derivatives, the Lagrangian and the field equations with
their numerical formulation, and four analytic solutions) is complete:
\texttt{superfluid.md}, \texttt{.tex}, \texttt{.pdf}, with the Rust
workspace \texttt{rust/} (program \texttt{sf\_part1}: VERDICT: ALL PASS
(76 of 76)) and the Python programs \texttt{python/sf\_part1.py}
(VERDICT: ALL PASS (33 of 33)) and
\texttt{python/sf\_part1\_figures.py}.

\textbf{Part II of the manuscript} (item 6 of the problem and the three
methods: the pair channels and the gauge bosons that mediate the
pairing, the methods HS, KS and CJT and their numerical engine, the
recursion of the potential to \(V^{(2)}_{\rm SF}\), the ground states,
the first excited states and the gaps, and the gauge bosons in the
superfluid) is complete, with the program \texttt{sf\_part2} (VERDICT:
ALL PASS (44 of 44)) and the Python programs
\texttt{python/sf\_part2.py} (VERDICT: ALL PASS (23 of 23)) and
\texttt{python/sf\_part2\_figures.py}.

\textbf{Part III of the manuscript} (the effective Lagrangians
\(\mathcal L^{\rm CJT}\), \(\mathcal L^{\rm HS}\),
\(\mathcal L^{\rm KS}\), their potential parts (the Mexican hat) and
kinetic parts, the vacua and the breaking of the symmetry, and
\(\mathrm{SFBSv}^X\)) is complete, with the program \texttt{sf\_part3}
(VERDICT: ALL PASS (31 of 31)) and the Python programs
\texttt{python/sf\_part3.py} (VERDICT: ALL PASS (10 of 10)) and
\texttt{python/sf\_part3\_figures.py}.

\textbf{Part IV of the manuscript} (the mass matrix of the gauge bosons
in every vacuum and its diagonalization, the theorem that the massless
gauge bosons are the generators that leave the condensate unchanged, the
surviving group \(SL(2,\mathbb C)\times O(2)\) and its \(U(1)\) with the
massless photon; the dynamics of the superfluid and of its gauge fields,
integrated with CVODE and IDA; two videos) is complete, with the program
\texttt{sf\_part4} (VERDICT: ALL PASS (120 of 120)), the integrators
\texttt{sfsundials} and their C interface \texttt{sfsundials\_c}, and
the Python programs \texttt{python/sf\_part4.py} (VERDICT: ALL PASS (20
of 20)) and \texttt{python/sf\_part4\_figures.py}. The Mathematica
notebook \texttt{superfluid.nb} is the next object of the plan.

\hypertarget{contents}{%
\section{3. Contents}\label{contents}}

\begin{longtable}[]{@{}
  >{\raggedright\arraybackslash}p{(\columnwidth - 4\tabcolsep) * \real{0.3333}}
  >{\raggedright\arraybackslash}p{(\columnwidth - 4\tabcolsep) * \real{0.3333}}
  >{\raggedright\arraybackslash}p{(\columnwidth - 4\tabcolsep) * \real{0.3333}}@{}}
\toprule\noalign{}
\begin{minipage}[b]{\linewidth}\raggedright
file or folder
\end{minipage} & \begin{minipage}[b]{\linewidth}\raggedright
what it is
\end{minipage} & \begin{minipage}[b]{\linewidth}\raggedright
state
\end{minipage} \\
\midrule\noalign{}
\endhead
\bottomrule\noalign{}
\endlastfoot
\texttt{README.md}, \texttt{.tex}, \texttt{.pdf} & this list of contents
& written \\
\texttt{provenance/P0\_task\_record.md}, \texttt{.tex}, \texttt{.pdf} &
the task (paraphrased), how the author's rules are kept, the record of a
false statement and its correction, the decisions of interpretation, the
plan, the log of errors and changes & written \\
\texttt{provenance/P1\_master\_symbols.md}, \texttt{.tex}, \texttt{.pdf}
& the MASTER provenance of the symbols of PROJECT spacetimes for this
solution: the convention of the symbols (suffix LL: taken unchanged from
the earlier work on the super-loop; suffix SF: introduced by this
solution), the seven registries of the project read by a program, the
identifications and remedies file by file, the collisions and their
resolution, the machine checks, and the table of all symbols & written;
grows with the solution \\
\texttt{provenance/symbols\_sf.csv} & the registry of the symbols of the
solution, written by \texttt{code/scripts/write\_symbol\_registry.py}
from \texttt{code/scripts/symbol\_entries.py} & written \\
\texttt{superfluid.md}, \texttt{.tex}, \texttt{.pdf} & the manuscript of
the solution & parts I to IV written \\
\texttt{superfluid.nb} & the Mathematica notebook of the whole solution,
with its provenance & planned \\
\texttt{provenance/} (further documents) & the provenance of the code,
the computed values and the checks, and the provenance of the notebook &
planned \\
\texttt{code/scripts/} & the document pipeline:
\texttt{assemble\_docs.py} (writes a document from its source in
\texttt{docs\_src/} and the generated fragments),
\texttt{build\_docs.sh} (the \texttt{.tex} and the \texttt{.pdf} of
every document, with a check of the layout of every page),
\texttt{doc\_header.tex}, \texttt{keep\_tables.py},
\texttt{check\_layout.py}, \texttt{check\_pages.py}, and
\texttt{check\_self\_contained.py} (does a document send its reader
elsewhere?); the registry of the symbols: \texttt{symbol\_entries.py}
(the entries), \texttt{write\_symbol\_registry.py} (reads the seven
registries of the project, writes the registry of the solution and the
tables of the MASTER provenance) and \texttt{check\_symbols.py} (the
checks of the symbols) & written \\
\texttt{docs\_src/} & the hand-edited sources of the documents &
written \\
\texttt{rust/} & the Rust workspace: the library \texttt{sfcore}
(bi-quaternions and bi-octonions, Dirac matrices, automatic
differentiation to second order, Fock-space operators, eigenvalues of
block-tridiagonal matrices, quadrature, reports; for part II the uniform
quasi-free states of the three methods, the linear response, the
eigenvalues of general complex matrices and the exchange of gauge
bosons; for part IV the response at complex frequencies), the program
\texttt{sf\_part1} (the checks and analytic solutions of part I), the
program \texttt{sf\_part2} (part II), the program \texttt{sf\_part3}
(part III), the program \texttt{sf\_part4} (part IV), the integrators
CVODE and IDA of the vendored pure-Rust SUNDIALS 7.8.0 with a plain
interface (\texttt{sfsundials}), and their C interface for Python, the
shared library \texttt{libsfsundials\_c} (\texttt{sfsundials\_c}) &
parts I to IV written \\
\texttt{python/} & \texttt{sfcommon.py} (common definitions),
\texttt{sfsundials.py} (CVODE and IDA through the C interface),
\texttt{sf\_part1.py}, \texttt{sf\_part2.py}, \texttt{sf\_part3.py},
\texttt{sf\_part4.py} (the second, independent computations of parts I
to IV), \texttt{sf\_part1\_figures.py}, \texttt{sf\_part2\_figures.py},
\texttt{sf\_part3\_figures.py}, \texttt{sf\_part4\_figures.py} (the
figures, and the videos of part IV) & parts I to IV written \\
\texttt{wolfram/} & the Wolfram-language parts of the notebook: in
\texttt{wolfram/tools/} the builder of the notebook from its source
(\texttt{build\_notebook.wls}), the reader of formulas that turns the
TeX of the texts into two-dimensional formulas
(\texttt{sf\_formula.wl}), and the evaluation of the notebook in the
front end of Wolfram (\texttt{guirun.wl}); the sources of the notebook
follow with object O8 & tools written; sources planned \\
\texttt{data/} & the reports of the checks
(\texttt{hck\_part1\_checks.txt},
\texttt{hck\_part1\_checks\_python.txt},
\texttt{hck\_part2\_checks.txt},
\texttt{hck\_part2\_checks\_python.txt},
\texttt{hck\_part3\_checks.txt},
\texttt{hck\_part3\_checks\_python.txt},
\texttt{hck\_part4\_checks.txt},
\texttt{hck\_part4\_checks\_python.txt},
\texttt{hck\_symbol\_checks.txt},
\texttt{hck\_self\_contained\_checks.txt}) and the data of parts I to IV
(\texttt{part1\_*.csv}, \texttt{part2\_*.csv}, \texttt{part3\_*.csv},
\texttt{part4\_*.csv}) & parts I to IV written \\
\texttt{figures/}, \texttt{videos/} & the figures
(\texttt{fig\_part1\_*}, \texttt{fig\_part2\_*}, \texttt{fig\_part3\_*},
\texttt{fig\_part4\_*}, each \texttt{.pdf} and \texttt{.png}) and the
videos (\texttt{video\_part4\_v1\_condensate\_growth.mp4},
\texttt{video\_part4\_v2\_gauge\_fields.mp4}) & written \\
\texttt{code/scripts/write\_part1\_fragments.py},
\texttt{write\_part2\_fragments.py},
\texttt{write\_part3\_fragments.py}, \texttt{write\_part4\_fragments.py}
& turn the reports and values of parts I to IV into the tables and
numbers of the manuscript & written \\
\end{longtable}

\hypertarget{how-to-build-the-documents}{%
\section{4. How to build the
documents}\label{how-to-build-the-documents}}

The programs are run first, then the documents are built, from the
folder \texttt{Revision\_HS\_CJT\_and\_Kohn-Sham}:

\begin{Shaded}
\begin{Highlighting}[]
\BuiltInTok{cd} \StringTok{"/Users/nsh/Developer/github/spacetimes/Revision\_HS\_CJT\_and\_Kohn{-}Sham"}
\KeywordTok{(}\BuiltInTok{cd}\NormalTok{ rust }\KeywordTok{\&\&} \ExtensionTok{cargo}\NormalTok{ build }\AttributeTok{{-}{-}release} \KeywordTok{\&\&} \ExtensionTok{./target/release/sf\_part1}\KeywordTok{)}   \CommentTok{\# part I: checks and data}
\ExtensionTok{python3}\NormalTok{ python/sf\_part1.py                    }\CommentTok{\# part I: the Python computation of the checks}
\ExtensionTok{python3}\NormalTok{ python/sf\_part1\_figures.py            }\CommentTok{\# part I: the figures}
\ExtensionTok{python3}\NormalTok{ code/scripts/write\_part1\_fragments.py }\CommentTok{\# the tables and numbers of part I for the manuscript}
\KeywordTok{(}\BuiltInTok{cd}\NormalTok{ rust }\KeywordTok{\&\&} \ExtensionTok{./target/release/sf\_part2}\KeywordTok{)}        \CommentTok{\# part II: checks and data (about nine minutes)}
\ExtensionTok{python3}\NormalTok{ python/sf\_part2.py                    }\CommentTok{\# part II: the Python computation (after sf\_part2)}
\ExtensionTok{python3}\NormalTok{ python/sf\_part2\_figures.py            }\CommentTok{\# part II: the figures}
\ExtensionTok{python3}\NormalTok{ code/scripts/write\_part2\_fragments.py }\CommentTok{\# the tables and numbers of part II for the manuscript}
\KeywordTok{(}\BuiltInTok{cd}\NormalTok{ rust }\KeywordTok{\&\&} \ExtensionTok{./target/release/sf\_part3}\KeywordTok{)}        \CommentTok{\# part III: checks and data (after sf\_part2)}
\ExtensionTok{python3}\NormalTok{ python/sf\_part3.py                    }\CommentTok{\# part III: the Python computation}
\ExtensionTok{python3}\NormalTok{ python/sf\_part3\_figures.py            }\CommentTok{\# part III: the figures}
\ExtensionTok{python3}\NormalTok{ code/scripts/write\_part3\_fragments.py }\CommentTok{\# the tables and numbers of part III for the manuscript}
\KeywordTok{(}\BuiltInTok{cd}\NormalTok{ rust }\KeywordTok{\&\&} \ExtensionTok{./target/release/sf\_part4}\KeywordTok{)}        \CommentTok{\# part IV: checks and data (after sf\_part3; about twenty minutes)}
\ExtensionTok{python3}\NormalTok{ python/sf\_part4.py                    }\CommentTok{\# part IV: the Python computation (uses the library of cargo build)}
\ExtensionTok{python3}\NormalTok{ python/sf\_part4\_figures.py            }\CommentTok{\# part IV: the figures and the two videos (needs ffmpeg)}
\ExtensionTok{python3}\NormalTok{ code/scripts/write\_part4\_fragments.py }\CommentTok{\# the tables and numbers of part IV for the manuscript}
\ExtensionTok{python3}\NormalTok{ code/scripts/write\_symbol\_registry.py }\CommentTok{\# the registry of the symbols and its tables}
\ExtensionTok{python3}\NormalTok{ code/scripts/check\_symbols.py         }\CommentTok{\# writes data/hck\_symbol\_checks.txt}
\ExtensionTok{python3}\NormalTok{ code/scripts/assemble\_docs.py         }\CommentTok{\# docs\_src/*.src.md {-}\textgreater{} the .md documents}
\FunctionTok{bash}\NormalTok{ code/scripts/build\_docs.sh               }\CommentTok{\# .md {-}\textgreater{} .tex {-}\textgreater{} .pdf, with the layout check}
\ExtensionTok{python3}\NormalTok{ code/scripts/check\_self\_contained.py  }\CommentTok{\# writes data/hck\_self\_contained\_checks.txt}
\end{Highlighting}
\end{Shaded}

The Rust programs need the Rust toolchain (\texttt{cargo}; no crate is
downloaded: the integrators are the pure-Rust SUNDIALS 7.8.0 vendored in
PROJECT spacetimes as \texttt{vendor/sundials\_rs-7.8.0}); the Python
programs need Python 3 with numpy, scipy, sympy and matplotlib, and the
videos the program ffmpeg. \texttt{assemble\_docs.py} writes each
document from its source in \texttt{docs\_src/}; \texttt{build\_docs.sh}
needs pandoc (version 2.11 or later), a TeX distribution with XeLaTeX
and latexmk, the DejaVu fonts and the font Latin Modern Math;
\texttt{check\_self\_contained.py} checks that no document refers its
reader to another document.

\end{document}
